\chapter{Mutation Testing}

\section{Limiti dei criteri di adequacy precedenti}

I criteri di adeguatezza visti in precedenza permettono di stimare quanto una test suite esplori il \textbf{System Under Test} (SUT), ma non misurano direttamente la sua capacità di rilevare difetti.

La \textbf{functional adequacy} si basa su specifiche, requisiti o comportamenti osservabili. Il suo limite principale è che non riesce a scoprire errori dovuti a dettagli implementativi introdotti dal programmatore.

La \textbf{structural adequacy}, invece, si basa sul codice e misura l'esplorazione di statement, branch o condizioni. Tuttavia, anche una coverage molto alta non garantisce che i test siano realmente efficaci nel trovare difetti.

\begin{notaslide}
I criteri funzionali e strutturali danno solo un'idea di quanta parte del SUT sia stata esplorata, ma non sono indici diretti dell'effettiva capacità di \emph{fault detection} della test suite.
\end{notaslide}

\begin{notaprof}
L'assunzione secondo cui esplorare tutto il codice faccia trovare i bug è solo una correlazione empirica. Se si raggiunge il 100\% di coverage, ma le asserzioni non controllano il reale stato del sistema, i test possono comunque passare ignorando errori logici.
\end{notaprof}

\begin{notaesame}
Alta coverage non significa automaticamente test efficaci. La coverage misura l'esplorazione del codice, non la forza degli oracle né la capacità della test suite di rilevare difetti.
\end{notaesame}

\section{Idea del Mutation Testing}

Il \textbf{Mutation Testing} è un approccio basato sull'inserimento deliberato di \textbf{mutazioni}, cioè alterazioni artificiali, all'interno del codice del SUT.

L'obiettivo non è misurare semplicemente quanto codice viene eseguito, ma stimare la bontà di una test suite osservando se essa riesce ad accorgersi di difetti artificialmente introdotti.

In questo senso, il Mutation Testing misura:

\begin{itemize}
    \item la potenziale capacità della test suite di rilevare errori;
    \item la robustezza della test suite rispetto a modifiche ed evoluzioni future del codice;
    \item la capacità degli oracle di distinguere il comportamento originale da quello alterato.
\end{itemize}

\begin{notaesame}
Il Mutation Testing valuta la qualità di una test suite esistente introducendo piccoli difetti artificiali nel SUT e verificando se i test riescono a rilevarli.
\end{notaesame}

\section{Concetti fondamentali}

Nel Mutation Testing si usano alcuni concetti ricorrenti.

\begin{description}
    \item[\(P\) -- SUT originale:] è il programma di partenza, assunto per ipotesi come corretto.

    \item[\(M\) -- Mutante:] è una variante del SUT in cui è stata introdotta una sola alterazione artificiale.

    \item[\(T\) -- Test set:] è l'insieme dei test usati per valutare il programma originale e i mutanti.

    \item[\(D\) -- Mutanti killed o distinguished:] sono i mutanti rivelati dai test, cioè quelli per cui almeno un test produce un comportamento diverso rispetto al programma originale.

    \item[\(L\) -- Mutanti live:] sono i mutanti vivi, cioè quelli che nessun test è riuscito a distinguere dal programma originale.

    \item[\(E\) -- Mutanti equivalenti:] sono mutanti sintatticamente diversi dal programma originale, ma semanticamente equivalenti, quindi indistinguibili da qualunque test.
\end{description}

\begin{notaslide}
Il test set \(T\) può essere visto come insieme di tuple di input oppure come insieme di \emph{test program} dotati di \texttt{assert}. Nel secondo caso, i test devono passare sul programma originale \(P\) e fallire sul mutante \(M\), se la mutazione è osservabile.
\end{notaslide}

\begin{notaprof}
Ogni mutante dovrebbe contenere una sola alterazione. Se un mutante contenesse più modifiche, un errore potrebbe mascherarne un altro e rendere difficile capire quale modifica sia stata effettivamente rilevata dai test.
\end{notaprof}

\section{Processo di Mutation Testing}

Il processo di Mutation Testing può essere descritto come una sequenza di passi.

\begin{enumerate}
    \item Si esegue il test set \(T\) sul programma originale \(P\). I test devono passare, perché \(P\) viene assunto come comportamento di riferimento.

    \item Si generano i mutanti \(M\), ciascuno ottenuto introducendo una piccola alterazione artificiale nel SUT.

    \item I mutanti generati vengono inseriti nell'insieme dei mutanti vivi \(L\).

    \item Per ogni mutante \(M\), si eseguono i test del test set \(T\).

    \item Per ogni test \(t\), si confronta il comportamento del programma originale con quello del mutante:
    \[
        P(t) = M(t)
    \]

    \item Se almeno un test produce un comportamento diverso:
    \[
        P(t) \neq M(t)
    \]
    allora il mutante viene considerato \textbf{killed} o \textbf{distinguished} e viene spostato nell'insieme \(D\).

    \item Se invece nessun test riesce a distinguere il mutante dal programma originale, il mutante resta vivo.

    \item Al termine del processo, i mutanti vivi vengono analizzati per individuare eventuali mutanti equivalenti \(E\).

    \item Infine, si calcola il \textbf{Mutation Score}.
\end{enumerate}

\begin{notaslide}
Nel diagramma del processo, le linee continue rappresentano il flusso di controllo, mentre le linee tratteggiate rappresentano il trasferimento di dati o artefatti.
\end{notaslide}

\begin{notaesame}
Un mutante viene ucciso quando almeno un test riesce a distinguere il comportamento del mutante dal comportamento del programma originale.
\end{notaesame}

\section{Mutation Score}

Il \textbf{Mutation Score} misura la capacità della test suite di uccidere i mutanti generati.

Una prima formulazione è:

\[
MS = \frac{|D|}{|L| + |D|}
\]

dove:

\begin{itemize}
    \item \(|D|\) è il numero di mutanti uccisi;
    \item \(|L|\) è il numero di mutanti vivi non equivalenti.
\end{itemize}

Una formulazione alternativa, che tiene conto dei mutanti equivalenti, è:

\[
MS = \frac{|D|}{|M| - |E|}
\]

dove:

\begin{itemize}
    \item \(|M|\) è il numero totale di mutanti generati;
    \item \(|E|\) è il numero di mutanti equivalenti;
    \item \(|M| - |E|\) rappresenta il numero di mutanti effettivamente rilevabili.
\end{itemize}

L'interpretazione è la seguente:

\begin{itemize}
    \item \textbf{Mutation Score = 1}: la test suite ha ucciso tutti i mutanti non equivalenti ed è considerata adeguata rispetto al criterio di mutation testing;
    \item \textbf{Mutation Score = 0}: nessun mutante è stato ucciso, quindi la test suite non è riuscita ad accorgersi di alcuna alterazione artificiale.
\end{itemize}

\begin{notaprof}
I mutanti equivalenti non possono essere rivelati da nessun test. Se venissero conteggiati nel denominatore, il Mutation Score risulterebbe artificialmente basso, penalizzando la test suite per un limite teorico del generatore di mutanti.
\end{notaprof}

\begin{notaesame}
I mutanti equivalenti devono essere esclusi dal calcolo del Mutation Score, perché nessun test potrebbe distinguerli dal programma originale.
\end{notaesame}

\section{Condizioni per uccidere un mutante}

Perché un mutante venga ucciso, devono essere soddisfatte contemporaneamente tre condizioni.

\begin{description}
    \item[C1 -- Raggiungibilità:] deve esistere un cammino di esecuzione che raggiunga lo statement mutato.

    \item[C2 -- Infezione dello stato:] dopo l'esecuzione dello statement mutato, lo stato interno del mutante deve differire dallo stato interno del programma originale.

    \item[C3 -- Propagazione dello stato:] la differenza di stato deve propagarsi fino a un punto osservabile, come l'output finale o una \texttt{assert}.
\end{description}

Se un test raggiunge il codice mutato ma la mutazione non altera lo stato, il mutante non viene ucciso. Se altera lo stato ma tale differenza viene sovrascritta prima dell'osservazione, il mutante continua comunque a sopravvivere.

\begin{notaesame}
Per uccidere un mutante non basta eseguire la riga mutata: la mutazione deve modificare lo stato e tale modifica deve propagarsi fino a un comportamento osservabile.
\end{notaesame}

\begin{notaesame}
Se non esiste alcun test nel dominio degli input capace di soddisfare contemporaneamente C1, C2 e C3, allora il mutante è equivalente.
\end{notaesame}

\section{Principi importanti}

\subsection{Competent Programmer Hypothesis}

La \textbf{Competent Programmer Hypothesis} afferma che un programmatore competente produce un programma non troppo distante dalla versione corretta.

L'idea alla base del Mutation Testing è che molti difetti reali siano piccoli errori locali, come operatori sbagliati, condizioni invertite o modifiche sintattiche minime. Per questo motivo, introdurre mutazioni artificiali piccole può aiutare a stimare la capacità della test suite di rilevare difetti realistici.

\begin{notaprof}
L'assunzione è che il programmatore abbia compreso l'algoritmo e che il programma scritto sia vicino a quello corretto, ma possa comunque contenere piccoli errori sintattici, logici o tipografici.
\end{notaprof}

\subsection{Indecidibilità dell'equivalenza}

Stabilire in modo generale e automatico se un mutante è semanticamente equivalente al programma originale è un problema \textbf{indecidibile}.

Per questo motivo, nella pratica, l'identificazione dei mutanti equivalenti richiede spesso analisi manuale.

\begin{notaesame}
Il tester ha un ruolo importante nell'analisi dei mutanti sopravvissuti: deve capire se sono dovuti a test deboli, oracle incompleti, codice non raggiunto oppure mutanti equivalenti.
\end{notaesame}

\subsection{Strong mutation e weak mutation}

Si distinguono due forme principali di mutation testing.

Nella \textbf{strong mutation}, un mutante è considerato ucciso solo se produce un comportamento osservabile esternamente diverso dal programma originale.

Nella \textbf{weak mutation}, invece, l'osservazione può fermarsi a stati intermedi interni, per esempio valori di variabili locali subito dopo lo statement mutato.

\begin{notaesame}
Un mutante può essere equivalente rispetto alla strong mutation, perché non cambia l'output osservabile, ma non equivalente rispetto alla weak mutation, perché produce uno stato interno diverso.
\end{notaesame}

\section{PIT / Pitest}

\textbf{PIT}, o \textbf{Pitest}, è un framework Java usato per automatizzare il Mutation Testing.

PIT si integra con il ciclo di build e può essere usato insieme a Maven e Surefire. Il tool modifica il bytecode, genera mutanti, esegue i test e produce un report sui mutanti uccisi e sopravvissuti.

\begin{notaslide}
PIT si integra nel progetto Maven tramite il plugin \texttt{pitest-maven}. La configurazione può essere inserita nel \texttt{pom.xml}, specificando le classi da mutare e i test da eseguire.
\end{notaslide}

Gli elementi principali della configurazione sono:

\begin{itemize}
    \item \texttt{targetClasses}: indica le classi del SUT da mutare;
    \item \texttt{targetTests}: indica le classi di test da usare;
    \item goal \texttt{mutationCoverage}: avvia l'analisi di mutation testing;
    \item directory \texttt{target/pit-reports}: contiene il report HTML generato.
\end{itemize}

Il goal Maven può essere eseguito con un comando del tipo:

\begin{verbatim}
mvn org.pitest:pitest-maven:mutationCoverage
\end{verbatim}

\begin{notaprof}
Nel report HTML di PIT è possibile aprire una classe e osservare, riga per riga, quante mutazioni sono state generate, quante sono state uccise e quante sono sopravvissute. Questo permette di capire quali aspetti del codice hanno ingannato la test suite.
\end{notaprof}

\section{Collegamento con il progetto}

Nel progetto, il Mutation Testing può essere usato per valutare la qualità reale dei test già scritti.

\begin{notaesame}
JaCoCo e PIT misurano aspetti diversi. JaCoCo misura l'esplorazione white-box del codice; PIT misura la capacità della test suite di rilevare difetti artificiali.
\end{notaesame}

Una possibile strategia è usare JaCoCo per capire quali parti del codice non sono state ancora esplorate e PIT per capire se i test già presenti contengono oracle sufficientemente forti.

Nel report di progetto è utile indicare:

\begin{itemize}
    \item il tool usato;
    \item il Mutation Score ottenuto;
    \item il numero di mutanti generati;
    \item il numero di mutanti uccisi;
    \item il numero di mutanti sopravvissuti;
    \item eventuali mutanti equivalenti o sospetti equivalenti;
    \item le cause principali dei mutanti sopravvissuti.
\end{itemize}

Se sopravvivono molti mutanti, bisogna analizzare il report HTML e distinguere tra diverse cause possibili:

\begin{itemize}
    \item test senza \texttt{assert};
    \item oracle deboli;
    \item codice eseguito ma non verificato;
    \item rami difficili da raggiungere;
    \item mutanti equivalenti;
    \item mutazioni in parti poco rilevanti del codice.
\end{itemize}

\begin{notaesame}
Non bisogna riportare solo il numero finale prodotto da PIT. Bisogna spiegare quali mutanti sono sopravvissuti e perché.
\end{notaesame}
